% answers_short.tex —— 2009 年·普通高等学校招生考试·江西卷·理科 · 简短答案速查
% 由 tools/make_short_answers.py 自 answers.tex 的 \gksolhead/\gkanswer 提取，
% 并与 2009/2009_数学真题_江西理_答案要点.md 逐题比对后生成，勿手改（改源 MD 后重跑）。

\begin{gknotice}
\textbf{使用说明}：本节只给最终答案便于快速核对，完整推导见解析版。答案为编写者独立推导，\textbf{非官方标准答案}。
\end{gknotice}

\gkpart{选择题}
\begin{center}\begin{tabular}{ccccccccccccc}
\toprule
题号 & 1 & 2 & 3 & 4 & 5 & 6 & 7 & 8 & 9 & 10 & 11 & 12 \\
\midrule
答案 & A & C & D & B & A & B & D & A & B & D & C & B \\
\bottomrule
\end{tabular}\end{center}

\gkpart{填空题}
\begin{itemize}[leftmargin=2.6em]
  \item[\textbf{13.}] $5$。
  \item[\textbf{14.}] $8$。
  \item[\textbf{15.}] $\sqrt2$。
  \item[\textbf{16.}] B、C。
\end{itemize}

\gkpart{解答题}
\begin{itemize}[leftmargin=2.6em]
  \item[\textbf{17.}] （1）减区间 $(-\infty,0)$、$(0,1)$，增区间 $(1,+\infty)$；（2）按 $k$ 分类，见解析。
  \item[\textbf{18.}] （1）分布列见解析；（2）$E\xi=15$。
  \item[\textbf{19.}] （1）$A=45^\circ$，$C=60^\circ$；（2）$a=2\sqrt2$，$c=2\sqrt3$。
  \item[\textbf{20.}] （1）证明见解析；（2）$\arccos\dfrac{\sqrt3}{3}$；（3）$\dfrac{10\sqrt6}{27}$。
  \item[\textbf{21.}] （1）$\dfrac{x^2}{2b^2}-\dfrac{y^2}{25b^2}=1$；（2）证明见解析，定点为 $(\pm5b,0)$。
  \item[\textbf{22.}] （1）$a_n=\dfrac{3^n-1}{3^n+1}$；（2）证明见解析，可取 $\lambda=\max\left\{a,\dfrac1a\right\}$。
\end{itemize}
